\documentclass[12pt, a4paper]{article}
\usepackage[T1]{fontenc}
\usepackage{lmodern}
\usepackage{amsmath, amsthm, amssymb, mathtools}
\usepackage[margin=1in]{geometry}
\usepackage{xr-hyper}
\usepackage[colorlinks=true,linkcolor=blue,citecolor=blue]{hyperref}
\usepackage{cleveref}
\usepackage{graphicx, subcaption}
\usepackage{natbib}
\usepackage{booktabs}
\usepackage{enumerate}

% Code references: scaled monospace with line-breaking
\newcommand{\code}[1]{\texttt{\small #1}}

\newtheorem{theorem}{Theorem}[section]
\newtheorem{lemma}[theorem]{Lemma}
\newtheorem{proposition}[theorem]{Proposition}
\newtheorem{corollary}[theorem]{Corollary}
\newtheorem{definition}{Definition}[section]
\newtheorem{remark}{Remark}[section]
\newtheorem{conjecture}[theorem]{Conjecture}
\newtheorem{derivation}[theorem]{Derivation}
\newtheorem{observation}[theorem]{Observation}
\numberwithin{equation}{section}

\newcommand{\dd}{\mathrm{d}}
\newcommand{\seriestitle}{Why Rotating Fluids Make Polygons}


\externaldocument[I-]{../paper-1-mathematics/main}
\externaldocument[II-]{../paper-2-physics/main}
\externaldocument[III-]{../paper-3-gravity/main}
\externaldocument[IV-]{../paper-4-field-theory/main}
\externaldocument[A-]{../paper-A-appendices/main}

\title{\seriestitle\\[6pt]\large Paper V: The $S^3$ Framework and the ADE Phase Transition}
\author{Gordon Speagle}
\date{}

\begin{document}
\maketitle

\begin{abstract}
The polygon vortex framework (Papers~I--IV) operates on two-dimensional
surfaces.  We lift it to the three-sphere~$S^3$ via the Hopf fibration,
showing that the 600-cell---the 120 vertices of the binary icosahedral
group~$I^*$ viewed as unit quaternions---is the natural UV configuration.
The scalar Green's function on~$S^3$ is derived from first principles
($G(\chi) = (1/4\pi^2)(\pi - \chi)\cot\chi$, verified by direct PDE
substitution), and its spectral decomposition reveals all nine~$I^*$
irreps with degeneracies~$d_\rho^2$ (the regular representation).
The Onsager selection on~$S^3$ and the Schur conservation theorem
(the non-abelian analogue of Lax integrability) protect the 600-cell
phase.  The Hopf projection $S^3 \to S^2$ recovers the icosahedron.
At $K = 0$---the topological transition from $S^3$ spatial geometry
to~$\mathbb{R}^3$---the dominant saddle changes from the icosahedron
($N = 12$, $A_5$, $E_8$) to the heptagon ($N = 7$, $\mathbb{Z}_7$, SM),
a first-order phase transition in the Ehrenfest sense.  The integer-spin
sector (120~dimensions) is the classical vortex physics visible
on~$S^2$; the half-integer sector (128~dimensions) is the $E_8$
spinor, invisible classically but revealed on~$S^3$.
\end{abstract}

\tableofcontents

%% ====================================================================
\section{Introduction}
\label{sec:intro}
%% ====================================================================

Papers~I--IV develop the vortex framework on two-dimensional surfaces:
the flat plane~$\mathbb{R}^2$, the sphere~$\mathbf{S}^2$, and the
hyperbolic plane~$\mathbf{H}^2$.  The stability threshold $N = 7$,
the Chern--Simons gauge theory, and the Standard Model all emerge from
the Havelock eigenvalue decomposition on these surfaces.

This paper lifts the framework to three dimensions.  The key geometry
is the three-sphere~$S^3$, which fibers over~$S^2$ via the Hopf
fibration $S^1 \to S^3 \to S^2$.  The three Thurston model geometries
relevant to the vortex universe are:
\begin{center}
\begin{tabular}{lcll}
\toprule
Curvature & Geometry & Vortex ground state & Gauge group \\
\midrule
$K > 0$ & $S^3$ (Hopf) & Icosahedron ($N = 12$, $A_5$) & $E_8$ \\
$K = 0$ & $\mathbb{R}^3$ (trivial) & Heptagon ($N = 7$, $\mathbb{Z}_7$) & SM \\
$K < 0$ & $\widetilde{\mathrm{SL}}(2,\mathbb{R})$ (Seifert) &
  Heptagon ($N = 7$, $\mathbb{Z}_7$) & SM \\
\bottomrule
\end{tabular}
\end{center}
The Euler class jumps: $e = 1$ (Hopf) $\to$ $e = 0$ (trivial) $\to$
$e = N/2$ (Seifert).  The transition is topological, not metric.

%% ====================================================================
\section{The scalar Green's function on $S^3$}
\label{sec:green}
%% ====================================================================

\subsection{Eigenvalues and eigenfunctions}

On the unit three-sphere, the scalar Laplacian has eigenvalues
$-\Delta f = l(l+2)\,f$ for $l = 0, 1, 2, \ldots$, with
degeneracy $(l+1)^2$.  The zonal eigenfunctions (depending only on
geodesic distance~$\chi$) are the Gegenbauer polynomials $C_l^1(\cos\chi)$,
also known as Chebyshev polynomials of the second kind.
They satisfy the recurrence $C_0^1 = 1$, $C_1^1 = 2x$,
$C_{l+1}^1 = 2x\,C_l^1 - C_{l-1}^1$,
with $C_l^1(1) = l + 1$ and $C_l^1(\cos\chi) = \sin((l{+}1)\chi)/\sin\chi$.

\subsection{Closed form and verification}

The Green's function $G(\chi)$ satisfying
$\Delta G = -\delta + 1/\mathrm{Vol}(S^3)$ (where $\mathrm{Vol}(S^3) = 2\pi^2$)
is constructed from the eigenfunction expansion using the $S^3$ addition
theorem and summed to a closed form.

\begin{theorem}[$S^3$ Green's function]\label{thm:green-s3}
\begin{equation}\label{eq:green-s3}
  G(\chi) = \frac{1}{4\pi^2}\,(\pi - \chi)\,\cot\chi.
\end{equation}
\end{theorem}

\begin{proof}
Direct substitution into the $S^3$ radial Laplacian
$\Delta f = f'' + 2\cot(\chi)\,f'$.  Let $h(\chi) = (\pi - \chi)\cot\chi$.
Then $h' = -\cot\chi - (\pi - \chi)\csc^2\chi$ and
$h'' = 2\csc^2\chi + 2(\pi - \chi)\csc^2\chi\,\cot\chi$.
Computing $\Delta h = h'' + 2\cot\chi\cdot h'$:
the $(\pi - \chi)$ terms cancel, leaving
$\Delta h = 2(\csc^2\chi - \cot^2\chi) = 2$.
Therefore $\Delta G = 2/(4\pi^2) = 1/(2\pi^2) = 1/\mathrm{Vol}(S^3)$.
\end{proof}

Properties: $G(0^+) \to +\infty$ (Coulomb singularity,
$\sim \pi/(4\pi^2\chi)$); $G(\pi/2) = 0$;
$G(\pi^-) = -1/(4\pi^2)$ (finite at the antipode).

\begin{remark}
The $S^3$ Laplacian identity $\Delta_{S^3}(-G) = -1/(2\pi^2)$ (constant)
is the three-dimensional analogue of the $S^2$ identity
$\Delta_{S^2}[-\ln\sin(d/2)] = 1/2$ that underlies the eigenvalue
pairing theorem (Theorem~I-\ref{I-thm:hessian-pairing}).
\end{remark}

%% ====================================================================
\section{The 600-cell: $I^*$ on $S^3$}
\label{sec:600cell}
%% ====================================================================

The binary icosahedral group $I^* \subset \mathrm{SU}(2)$ has 120
elements, which we represent as unit quaternions
$q = w + x\mathbf{i} + y\mathbf{j} + z\mathbf{k} \in S^3$.
These 120 quaternions are the vertices of the 600-cell, a regular
4-dimensional polytope with 120~vertices, 720~edges,
1200~triangular faces, and 600~tetrahedral cells.

The construction: 8~quaternion units ($\pm 1, \pm\mathbf{i},
\pm\mathbf{j}, \pm\mathbf{k}$), 16~half-integer quaternions
$(\pm 1 \pm \mathbf{i} \pm \mathbf{j} \pm \mathbf{k})/2$,
and 96~golden quaternions (even permutations of
$(0, \pm 1, \pm\varphi, \pm 1/\varphi)/2$ where
$\varphi = (1+\sqrt{5})/2$).

\subsection{Distance classes and conjugacy classes}

The geodesic distance between quaternions $q_1, q_2 \in S^3$ is
$\chi = \arccos(q_1 \cdot q_2)$ (the $\mathbb{R}^4$ inner product).
The 119~non-identity 600-cell vertices fall into exactly eight
distance classes from any reference vertex, with sizes
$\{12, 20, 12, 30, 12, 20, 12, 1\}$.

These eight distance classes are precisely the eight nontrivial
conjugacy classes of~$I^*$ (the ninth class is the identity).
The geometry of the 600-cell IS the group theory of~$I^*$.

%% ====================================================================
\section{Energy decomposition: the regular representation}
\label{sec:energy-decomp}
%% ====================================================================

The energy weight matrix $K_{jk} = -G(\chi_{jk})$ for $j \neq k$,
$K_{jj} = -\sum_{l \neq j} K_{jl}$ (zero-sum rows), gives the
decomposition of the interaction energy by~$I^*$ irreps.

\begin{theorem}[Regular representation]\label{thm:regular-rep}
The eigenvalue degeneracies of the energy weight matrix on the
600-cell are $d_\rho^2$ for each $I^*$ irrep~$\rho$:
$\{1, 4, 9, 16, 25, 36, 16, 4, 9\}$.
\end{theorem}

\noindent
This is the regular representation structure: each irrep~$\rho$ of
dimension~$d_\rho$ appears with multiplicity~$d_\rho$ in the
permutation representation of $I^*$ acting on itself.

The nine distinct eigenvalues correspond to the nine $I^*$ irreps,
giving the energy-weighted Havelock Casimir $T_\rho$ for \emph{all}
nine irreps---including the four half-integer-spin irreps
($\rho_1, \rho_3, \rho_5, \rho_7$) that are invisible on~$S^2$
(where the icosahedron permutation representation factors through
$A_5 = I^*/\mathbb{Z}_2$ and sees only integer-spin irreps).

\begin{remark}
On~$S^2$, the icosahedron's 12-vertex permutation representation
decomposes as $1 \oplus 3 \oplus 5 \oplus 3'$ (four $A_5$ irreps,
all integer-spin).  On~$S^3$, the 600-cell's 120-vertex regular
representation decomposes into all nine $I^*$ irreps with multiplicity
$d_\rho$.  The $S^3$ framework \emph{completes} the $S^2$ picture.
\end{remark}

%% ====================================================================
\section{Onsager selection and Schur conservation}
\label{sec:onsager-schur}
%% ====================================================================

\subsection{Onsager selection on $S^3$}

At negative temperature (the Onsager regime), the equilibrium
maximises the energy~$H$.  Among $I^*$-symmetric configurations
on~$S^3$, the 600-cell has the maximum number of pairwise interactions
(120~vertices, $\binom{120}{2} = 7140$ pairs) and is selected by
the Onsager principle as the ground state.

\subsection{Schur conservation theorem}

\begin{theorem}[Schur conservation]\label{thm:schur}
For any $I^*$-equivariant perturbation of the 600-cell, the
energy weight matrix decomposes into $I^*$ irrep blocks with
eigenvalues that vary continuously within each block.
The block structure (degeneracies~$d_\rho^2$) is topologically
protected.
\end{theorem}

\begin{proof}
The energy weight matrix commutes with the $I^*$ action
(by $I^*$-equivariance of the configuration).  By Schur's lemma,
its eigenspaces are $I^*$ irreps.  The eigenvalue within each
block varies continuously under smooth perturbation, but the
\emph{degeneracy} $d_\rho^2$ is a discrete invariant (determined
by the representation type) that cannot change without breaking
$I^*$ symmetry.
\end{proof}

This is the $S^3$ analogue of the CMS Lax conservation on
$\mathbb{R}^2/\mathbf{H}^2$ (Paper~III, \S3): the Lax spectrum
prevents polygon--BTZ transitions; the Schur block structure
prevents 600-cell--lower-symmetry transitions.

%% ====================================================================
\section{The Hopf projection: $S^3 \to S^2$}
\label{sec:hopf}
%% ====================================================================

The Hopf fibration $S^1 \to S^3 \to S^2$ has Euler class $e = 1$
(the first Chern class of $\mathcal{O}(1)$ on $\mathbb{CP}^1 = S^2$).
It connects the 600-cell on~$S^3$ to the icosahedron on~$S^2$.

Each icosahedron vertex has a stabiliser $\mathbb{Z}_{10} \subset I^*$
(the ten $I^*$ elements whose $\mathrm{SO}(3)$ projection fixes that
vertex).  The group quotient $I^*/\mathbb{Z}_{10}$ gives the
12~icosahedron vertices: $120/10 = 12$.

The bridge identity
(Theorem~I-\ref{I-thm:bridge}, eq.~\eqref{I-eq:bridge})
\[
  \lambda_j = 5K_1\, P_j\!\Bigl(\frac{1}{\sqrt{5}}\Bigr)
            + 5K_2\, P_j\!\Bigl(-\frac{1}{\sqrt{5}}\Bigr)
            + \frac{1}{4}(-1)^j
\]
is the Hopf descendant of the $S^3$ spectral decomposition:
the Gegenbauer eigenvalues on~$S^3$ project to the Legendre
eigenvalues on~$S^2$ through the Hopf map.

%% ====================================================================
\section{The $S^3$ Havelock decomposition}
\label{sec:s3-havelock}
%% ====================================================================

The $S^3$ Havelock Casimir for the $l$-th harmonic on the 600-cell is
\begin{equation}\label{eq:s3-havelock}
  T_l = \sum_{k \neq 0} w(\chi_{0k})
  \Bigl[1 - \frac{C_l^1(\cos\chi_{0k})}{l+1}\Bigr],
\end{equation}
where $w = -G$ is the interaction weight and $C_l^1/(l+1)$ is the
normalised zonal spherical function on~$S^3$.

The corresponding eigenvalue
$\lambda_l = C_1 - T_l = \sum_{k \neq 0} w(\chi_{0k})\,
C_l^1(\cos\chi_{0k})/(l+1)$
is the $S^3$ bridge formula (the Gegenbauer analogue of the
$S^2$ Legendre formula).

%% ====================================================================
\section{The phase transition at $K = 0$}
\label{sec:phase-transition}
%% ====================================================================

\begin{theorem}[First-order phase transition]\label{thm:phase-transition}
The dominant vortex configuration changes discontinuously at $K = 0$:
\begin{itemize}
\item $K > 0$ ($S^2$): Icosahedron ($N = 12$, $A_5$, $E_8$ via McKay).
\item $K = 0$ ($\mathbb{R}^2$): Heptagon ($N = 7$, $\mathbb{Z}_7$,
  $\mathrm{SU}(3) \times \mathrm{SU}(2) \times \mathrm{U}(1)$
  via Frobenius).
\end{itemize}
\end{theorem}

\emph{Thermodynamic framework.}\enspace
Microcanonical ensemble at fixed angular impulse~$L$,
negative temperature $\beta < 0$.
The order parameter $\sigma(K) = $ symmetry group of the dominant
saddle: $\sigma(K > 0) = A_5$ (non-abelian, $|G| = 60$),
$\sigma(K = 0) = \mathbb{Z}_7$ (abelian, $|G| = 7$).
The jump is discontinuous $\Rightarrow$ first-order by Ehrenfest.

\emph{The mechanism.}\enspace
On $S^2(R)$ for any $R > 0$, the icosahedron has higher energy
than the heptagon:
$\Delta E(R) = \Delta E_{\mathrm{unit}} + 45\ln R \to \infty$.
But polygon rings are \emph{all unstable} on~$S^2$
(Morse index $N - 5$ for $N \ge 8$), while the icosahedron is stable.
At $K = 0$, the topology changes: $S^2 \to \mathbb{R}^2$, and
Platonic configurations cease to exist as finite-energy objects.
The 236~coset generators ($E_8 \to \mathrm{SM}$, dimension $248 - 12$)
become massive at the transition.

%% ====================================================================
\section{The $120/128$ split: classical and quantum sectors}
\label{sec:split}
%% ====================================================================

The Coxeter decomposition
(Theorem~I-\ref{I-thm:coxeter}, eq.~\eqref{I-eq:coxeter})
gives $248 = 120 + 128$ under $I^*$:
\begin{itemize}
\item \textbf{Integer-spin sector} (120~dimensions,
  $\mathrm{SO}(16)$ adjoint):
  the five irreps $\rho_0, \rho_2, \rho_4, \rho_6, \rho_8$.
  These are visible on~$S^2$ as the eigenvalues of the icosahedron's
  interaction matrix.  This is the classical vortex physics.

\item \textbf{Half-integer-spin sector} (128~dimensions,
  $\mathrm{SO}(16)$ spinor):
  the four irreps $\rho_1, \rho_3, \rho_5, \rho_7$.
  These are invisible on~$S^2$ (the icosahedron permutation representation
  factors through $A_5 = I^*/\mathbb{Z}_2$) but revealed on~$S^3$
  (the 600-cell regular representation sees all nine irreps).
  This is the quantum/spinor sector.
\end{itemize}

The half-integer irreps have zero multiplicity in the icosahedron's
permutation representation because the $\mathbb{Z}_{10}$ stabiliser
character sum vanishes for half-integer-spin characters
(proved: the sum $\sum_{h \in \mathbb{Z}_{10}} \chi_\rho(h) = 0$
for all half-integer~$\rho$).

This $120/128$ split is the mathematical origin of the boson/fermion
distinction in the $E_8$ framework.

%% ====================================================================
\section{Conclusion}
\label{sec:conclusion}
%% ====================================================================

The $S^3$ framework completes the vortex universe:
\begin{center}
\small
\begin{tabular}{lll}
\toprule
Framework & Ground state & Gauge group \\
\midrule
$S^3$ (Paper V) & 600-cell ($I^*$, 120 vertices) & $E_8$ (UV) \\
$S^2$ (Papers I--IV) & Icosahedron ($A_5$, 12 vertices) & $E_8$ (via McKay) \\
$\mathbb{R}^2$ (Papers I--IV) & Heptagon ($\mathbb{Z}_7$, 7 vertices) & SM (via Frobenius) \\
\bottomrule
\end{tabular}
\end{center}
The transition from $S^3$ to $\mathbb{R}^3$ at $K = 0$ is topological
(a Thurston geometry change), producing the $E_8 \to \mathrm{SM}$
breaking without a Higgs mechanism.  Paper~VI uses this as the
foundation for the cosmological parameters: the 600-cell is the
inflationary ground state, the $K = 0$ transition is reheating,
and $\Lambda \propto K$ is the residual curvature.

All results in this paper are backed by 77 computational tests
in the companion code repository
(\texttt{proofs/s3\_vortex\_dynamics.py},
 \texttt{proofs/e8\_casimir\_bridge.py},
 \texttt{proofs/phase\_transition.py},
 \texttt{proofs/coupling\_constants.py}).

\bibliographystyle{plainnat}
\bibliography{../shared/refs}

\end{document}
